\SourceLeaf{63}
\begin{center}\emph{Ein Verfahren zum Finden primitiver Wurzeln.}\end{center}

\Article{73}
Die Methoden zum Finden primitiver Wurzeln beruhen größtenteils auf Probieren. Verbindet man das, was wir in Artikel 55 gezeigt haben, mit dem, was wir weiter unten zur Lösung der Kongruenz $x^n\equiv1$ angeben, so hat man fast alles, was direkte Methoden leisten können. Euler räumt in \emph{Opusc. Analyt. T.} I, \emph{p.} 152 ein, dass es äußerst schwierig scheine, diese Zahlen zu bestimmen, und dass ihre Natur zu den tiefsten Geheimnissen der Zahlen gehöre. Durch Probieren lassen sie sich aber wie folgt recht bequem bestimmen. Geübte werden die langwierige Rechnung mit verschiedenen besonderen Kunstgriffen abkürzen können; diese lernt man durch Übung viel schneller als durch Regeln.

$1^0$. Man wähle eine beliebige zu $p$ (so bezeichnen wir stets den Modul) teilerfremde Zahl $a$. Meist verkürzt es die Rechnung, eine möglichst kleine Zahl zu nehmen, \emph{zum Beispiel} $2$. Dann bestimme man ihre Periode (Artikel 46), \emph{das heißt} die kleinsten Reste ihrer Potenzen, bis man eine Potenz $a^t$ erreicht, deren kleinster Rest $1$ ist\HistoricalNote{*)}{Man sieht unmittelbar, dass die Potenzen selbst nicht bekannt sein müssen. Denn der kleinste Rest jeder Potenz lässt sich leicht aus dem kleinsten Rest der vorigen gewinnen.}. Ist nun $t=p-1$, so ist $a$ eine primitive Wurzel.

$2^0$. Ist dagegen $t<p-1$, so nehme man eine andere Zahl $b$, die in der Periode von $a$ nicht vorkommt, und bestimme ebenso ihre Periode. Bezeichnet man den Exponenten, zu dem $b$ gehört, mit $u$, so erkennt man leicht, dass $u$ weder gleich $t$ noch ein echter Teiler davon sein kann. In beiden Fällen wäre nämlich $b^t\equiv1$, was unmöglich ist, da die Periode von $a$ alle Zahlen umfasst, deren Potenz mit dem Exponenten $t$ der Eins kongruent ist (Artikel 53). Ist $u$ $=p-1$, so ist $b$ eine primitive Wurzel. Ist $u$ zwar nicht $=p-1$, aber doch ein Vielfaches von $t$, so kennen wir nun eine Zahl, die zu einem größeren Exponenten gehört, und sind unserem Ziel näher: eine Zahl zu finden, die zum \emph{größten} Exponenten gehört. Ist $u$ weder $=p-1$ noch ein Vielfaches von $t$, so können wir trotzdem eine Zahl finden, deren Exponent größer als $t,u$ ist, nämlich gleich dem kleinsten gemeinsamen Vielfachen von $t,u$. Dieses sei $=y$. Man zerlege $y$ in zwei teilerfremde Faktoren $m,n$, von denen einer $t$, der andere $u$ teilt\HistoricalNote{\ensuremath{\dagger})}{Wie dies möglich ist, ergibt sich leicht aus Artikel 18. Man zerlege $y$ in verschiedene Primzahlen oder Potenzen verschiedener Primzahlen. Jeder dieser Faktoren teilt eine der Zahlen $t,u$ oder auch beide. Man ordne jeden der Zahl $t$ oder der Zahl $u$ zu, je nachdem, welche er teilt. Teilt er beide, so ist die Zuordnung beliebig. Das Produkt der $t$ zugeordneten Faktoren sei $=m$, das Produkt der übrigen $=n$. Dann erkennt man leicht, dass $m$ die Zahl $t$, $n$ die Zahl $u$ teilt und $mn=y$ gilt.}. Man bilde nun die Potenz mit dem Exponenten $\frac{t}{m}$ von $a$, setze sie $\equiv A$, und die
\SourceLeaf{64}Potenz mit dem Exponenten $\frac{u}{n}$ von $b$, die man $\equiv B\,(\operatorname{mod.}p)$ setze. Das Produkt $AB$ gehört dann zum Exponenten $n$. Denn offenbar gehört $A$ zum Exponenten $m$ und $B$ zum Exponenten $n$. Also gehört das Produkt $AB$ zu $mn$, weil $m,n$ teilerfremd sind. Das lässt sich genauso beweisen wie in Artikel 55, II.

$3^0$. Ist nun $y=p-1$, so ist $AB$ eine primitive Wurzel. Andernfalls muss man wie zuvor eine andere Zahl nehmen, die in der Periode von $AB$ nicht vorkommt. Diese ist entweder eine primitive Wurzel oder gehört zu einem größeren Exponenten als $y$, oder man kann wenigstens mit ihrer Hilfe wie zuvor eine Zahl finden, die zu einem größeren Exponenten als $y$ gehört. Da die Zahlen, die man durch Wiederholen dieses Verfahrens erhält, zu stetig wachsenden Exponenten gehören, findet man schließlich eine Zahl, die zum \emph{größten} Exponenten gehört, \emph{das heißt} eine primitive Wurzel, \emph{wie verlangt.}

\Article{74}
Ein Beispiel macht diese Regeln deutlicher. Sei $p=73$, und gesucht sei eine primitive Wurzel dafür. Versuchen wir zuerst die Zahl $2$, deren Periode lautet:

\[\setlength{\arraycolsep}{2.5pt}\begin{array}{@{}rrrrrrrrrrl@{}}
1.&2.&4.&8.&16.&32.&64.&55.&37.&1&\text{usw.}\\
0.&1.&2.&3.&4.&5.&6.&7.&8.&9&\text{usw.}
\end{array}\]

Da schon die Potenz mit dem Exponenten $9$ der Eins kongruent ist, ist $2$ keine primitive Wurzel. Man versuche eine andere Zahl, die in der Periode von $2$ nicht vorkommt, \emph{zum Beispiel} $3$, mit der Periode

\[\setlength{\arraycolsep}{2.5pt}\begin{array}{@{}rrrrrrrrrrrrrl@{}}
1.&3.&9.&27.&8.&24.&72.&70.&64.&46.&65.&49.&1&\text{usw.}\\
0.&1.&2.&3.&4.&5.&6.&7.&8.&9.&10.&11.&12&\text{usw.}
\end{array}\]

Also ist auch $3$ keine primitive Wurzel. Das kleinste gemeinsame Vielfache der Exponenten, zu denen $2,3$ gehören, \emph{also} der Zahlen $9,12$, ist aber $36$. Nach den Regeln des vorigen Artikels zerlegt man es in die Faktoren $9$ und $4$. Man muss also $2$ zur Potenz mit dem Exponenten $\frac{9}{9}$ erheben, \emph{das heißt} die Zahl $2$ selbst beibehalten, dagegen $3$ zur Potenz mit dem Exponenten $3$ erheben. Ihr Produkt ist $54$ und gehört somit zum Exponenten $36$. Berechnet man schließlich die Periode von $54$ und versucht erneut eine darin nicht enthaltene Zahl, \emph{zum Beispiel} $5$, so findet man, dass diese eine primitive Wurzel ist.

\SourceLeaf{65}
\begin{center}\emph{Verschiedene Sätze über Perioden und primitive Wurzeln.}\end{center}

\Article{75}
Bevor wir diesen Gegenstand verlassen, geben wir einige Sätze an, die wegen ihrer Einfachheit Beachtung verdienen dürften.

\emph{Das Produkt aller Glieder der Periode einer beliebigen Zahl ist $\equiv1$, wenn ihre Anzahl, also der Exponent, zu dem die Zahl gehört, ungerade ist, und $\equiv-1$, wenn dieser Exponent gerade ist.}

\emph{Beispiel.} Für den Modul $13$ besteht die Periode der Zahl $5$ aus den Gliedern $1,5,12,8$, deren Produkt $480\equiv-1\,(\operatorname{mod.}13)$ ist.

Nach demselben Modul besteht die Periode der Zahl $3$ aus den Gliedern $1,3,9$, deren Produkt $27\equiv1\,(\operatorname{mod.}13)$ ist.

\emph{Beweis.} Der Exponent, zu dem die Zahl gehört, sei $t$, ihr Index $\frac{p-1}{t}$. Bei geeigneter Wahl der Basis ist dies stets möglich (Artikel 71). Dann ist der Index des Produkts aller Periodenglieder

\[\equiv(1+2+3+\text{usw.}+t-1)\frac{p-1}{t}=\frac{(t-1)(p-1)}{2}\]

\emph{also} $\equiv0\,(\operatorname{mod.}p-1)$, wenn $t$ ungerade ist, und $\equiv\frac{p-1}{2}$, wenn $t$ gerade ist. Im ersten Fall ist daher das Produkt $\equiv1\,(\operatorname{mod.}p)$, im zweiten $\equiv-1\,(\operatorname{mod.}p)$ (Artikel 62). \mbox{\emph{Q. E. D.}}

\Article{76}
Ist die Zahl im vorigen Satz eine primitive Wurzel, so umfasst ihre Periode sämtliche Zahlen $1,2,3\mathpunct{.\,.\,.\,.\,.\,.}\,p-1$, deren Produkt daher stets $\equiv-1$ ist (denn $p-1$ ist stets gerade, mit der einzigen Ausnahme $p=2$, in der $-1$ und $+1$ gleichwertig sind). Diesen schönen Satz, den man gewöhnlich so ausspricht: \emph{Das Produkt aller Zahlen unterhalb einer gegebenen Primzahl, um eins vermehrt, ist durch diese Primzahl teilbar}, hat Waring zuerst angegeben und Wilson, Esq., zugeschrieben, \emph{Meditt. algebr. Ed.} 3, \emph{p.} 380. Doch keiner von beiden konnte ihn beweisen. Waring räumt ein, dass der Beweis umso schwieriger scheine, als sich keine \emph{Bezeichnung} erfinden lasse, die eine Primzahl ausdrücke. --- Unserer Ansicht nach sollte man solche Wahrheiten aber eher aus Begriffen als aus Bezeichnungen gewinnen. Später gab La Grange einen Beweis, \emph{Nouv. M\'em. de l'Ac. de Berlin,} 1771. Er beruht auf der
\SourceLeaf{66}Betrachtung der Koeffizienten, die beim Ausmultiplizieren des Produkts

\[x+1.\,x+2.\,x+3\mathpunct{.\,.\,.\,.}\,x+p-1\]

\noindent entstehen. Setzt man nämlich dieses Produkt

\[\equiv x^{p-1}+Ax^{p-2}+Bx^{p-3}+\text{usw.}+Mx+N\]

\noindent so sind die Koeffizienten $A,B$ usw. $M$ durch $p$ teilbar; $N$ ist dagegen $=1.2.3\mathpunct{.\,.\,.\,.}\,p-1$. Für $x=1$ ist das Produkt nun durch $p$ teilbar, dann aber $\equiv1+N\,(\operatorname{mod.}p)$. Also muss $1+N$ durch $p$ teilbar sein.

Schließlich gab Euler in \emph{Opusc. analyt. T.} I, \emph{p.} 329 einen Beweis, der mit dem hier vorgetragenen übereinstimmt. Wenn solche Männer diesen Satz ihrer Überlegungen wert fanden, hoffen wir, dass man uns einen weiteren Beweis nicht verübeln wird.

\Article{77}
Ist nach dem Modul $p$ das Produkt zweier Zahlen $a,b$ der Eins kongruent, so nennen wir die Zahlen $a,b$ mit Euler \emph{Partner}. Nach dem vorigen Abschnitt hat dann jede positive Zahl unter $p$ genau einen positiven Partner unter $p$. Unter den Zahlen $1,2,3\mathpunct{.\,.\,.\,.}\,p-1$ sind, wie man leicht zeigt, nur $1$ und $p-1$ ihre eigenen Partner. Denn solche Zahlen sind Wurzeln der Kongruenz $xx\equiv1$. Da diese zweiten Grades ist, kann sie nicht mehr als zwei Wurzeln haben, \emph{also} keine anderen als $1$ und $p-1$. Lässt man diese beiden weg, so sind die übrigen Zahlen $2.3\mathpunct{.\,.\,.\,.}\,p-2$ paarweise Partner. Ihr Produkt ist daher $\equiv1$, also das Produkt aller Zahlen $1,2,3\mathpunct{.\,.\,.}\,p-1$ $\equiv p-1$ oder $\equiv-1$. \mbox{\emph{Q. E. D.}}

\emph{Zum Beispiel} bilden für $p=13$ die Zahlen $2,3,4\mathpunct{.\,.\,.\,.}\,11$ folgende Paare: $2$ mit $7$; $3$ mit $9$; $4$ mit $10$; $5$ mit $8$; $6$ mit $11$. Denn $2.7\equiv1$; $3.9\equiv1$ usw. Damit ist $2.3.4\mathpunct{.\,.\,.\,.}\,11\equiv1$, also $1.2.3\mathpunct{.\,.\,.\,.}\,12\equiv-1$.

\Needspace{9\baselineskip}
\Article{78}
Man kann den Wilsonschen Satz aber allgemeiner so formulieren: \emph{Das Produkt aller Zahlen, die kleiner als eine beliebige gegebene Zahl $A$ und zu ihr teilerfremd sind, ist nach $A$ der negativ oder positiv genommenen Eins kongruent.} Negativ nimmt man die Eins, wenn $A$ von der Form $p^m$ oder $2p^m$ ist, wobei $p$ eine
\SourceLeaf{67}von $2$ verschiedene Primzahl bezeichnet, und außerdem für $A=4$; in allen übrigen Fällen positiv. Der von Wilson angegebene Satz ist im ersten Fall enthalten. --- \emph{Zum Beispiel} ist für $A=15$ das Produkt der Zahlen $1,2,4,7,8,11,13,14$ $\equiv1\,(\operatorname{mod.}15)$. Der Kürze wegen fügen wir keinen Beweis hinzu. Wir bemerken nur, dass er wie im vorigen Artikel geführt werden kann, außer dass die Kongruenz $xx\equiv1$ mehr als zwei Wurzeln haben kann, was besondere Überlegungen erfordert. Auch aus der Betrachtung der Indizes könnte man einen Beweis gewinnen, ähnlich wie in Artikel 75, wenn man das hinzunimmt, was wir gleich über nichtprime Moduln vortragen werden.

\Article{79}
Wir kehren zur Aufzählung weiterer Sätze zurück (Artikel 75).

\emph{Die Summe aller Glieder der Periode einer beliebigen Zahl ist $\equiv0$}, wie im Beispiel aus Artikel 75: $1+5+12+8=26\equiv0\,(\operatorname{mod.}13)$.

\emph{Beweis.} Die Zahl, deren Periode wir betrachten, sei $=a$, und der Exponent, zu dem sie gehört, $=t$. Dann ist die Summe sämtlicher Periodenglieder
\[\equiv1+a+aa+a^3+\text{usw.}+a^{t-1}\equiv\frac{a^t-1}{a-1}\,(\operatorname{mod.}p)\]
Nun ist aber $a^t-1\equiv0$. Daher ist diese Summe stets $\equiv0$ (Artikel 22), außer wenn $a-1$ durch $p$ teilbar ist, also $a\equiv1$. Diesen Fall müssen wir somit ausnehmen, wenn wir auch ein einziges Glied eine \emph{Periode} nennen wollen.

\Article{80}
\emph{Das Produkt aller primitiven Wurzeln ist $\equiv1$}, mit der einzigen Ausnahme $p=3$; dann gibt es nur eine primitive Wurzel, nämlich $2$.

\emph{Beweis.} Nimmt man eine beliebige primitive Wurzel als Basis, so sind die Indizes aller primitiven Wurzeln die zu $p-1$ teilerfremden Zahlen, die kleiner als diese Zahl sind. Deren Summe, \emph{also} der Index des Produkts aller primitiven Wurzeln, ist $\equiv0\,(\operatorname{mod.}p-1)$, und somit das Produkt $\equiv1\,(\operatorname{mod.}p)$. Denn ist $k$ zu $p-1$ teilerfremd, so offenbar auch $p-1-k$ zu $p-1$. Also bilden je zwei zu $p-1$ teilerfremde Zahlen eine durch $p-1$ teilbare Summe. (Dabei kann $k$ niemals gleich $p-1-k$ sein, außer für $p-1=2$, also $p=3$, welchen Fall wir ausgenommen haben. In allen übrigen Fällen ist $\frac{p-1}{2}$ offenbar zu $p-1$ nicht teilerfremd.)

\SourceLeaf{68}
\Article{81}
\emph{Die Summe aller primitiven Wurzeln ist entweder} $\equiv0$ (wenn $p-1$ durch eine \emph{Quadratzahl} teilbar ist) \emph{oder $\equiv\pm1\,(\mathit{mod.}p)$} (wenn $p-1$ ein Produkt verschiedener Primzahlen ist; bei gerader Anzahl nimmt man das positive, bei ungerader das negative Vorzeichen).

\emph{Beispiele.} $1^0$ Für $p=13$ sind die primitiven Wurzeln $2,6,7,11$, deren Summe $26\equiv0\,(\operatorname{mod.}13)$ ist.

$2^0$ Für $p=11$ sind die primitiven Wurzeln $2,6,7,8$, deren Summe $23\equiv+1\,(\operatorname{mod.}11)$ ist.

$3^0$ Für $p=31$ sind die primitiven Wurzeln $3,11,12,13,17.\,21,22,24$, deren Summe $123\equiv-1\,(\operatorname{mod.}31)$ ist.

\emph{Beweis.} Oben haben wir gezeigt (Artikel 55, II): Ist $p-1$ $=a^{\alpha}b^{\beta}c^{\gamma}$ usw. (wobei $a,b,c$ usw. verschiedene Primzahlen bezeichnen), und sind $A,B,C$ usw. beliebige Zahlen, die zu den jeweiligen Exponenten $a^{\alpha},b^{\beta},c^{\gamma}$ usw. gehören, so stellen sämtliche Produkte $ABC$ usw. primitive Wurzeln dar. Leicht lässt sich auch zeigen, dass sich jede primitive Wurzel durch ein solches Produkt darstellen lässt, und zwar auf genau eine Weise\HistoricalNote{*)}{Man bestimme die Zahlen $\mathfrak{a},\mathfrak{b},\mathfrak{c}$ usw. so, dass $\mathfrak{a}\equiv1\,(\operatorname{mod.}a^{\alpha})$ und $\equiv0\,(\operatorname{mod.}b^{\beta}c^{\gamma}\text{usw.})$; $\mathfrak{b}\equiv1\,(\operatorname{mod.}b^{\beta})$ und $\equiv0\,(\operatorname{mod.}a^{\alpha}c^{\gamma}\text{usw.})$ usw. gilt (siehe Artikel 32). Dann wird $\mathfrak{a}+\mathfrak{b}+\mathfrak{c}+\text{usw.}\equiv1\,(\operatorname{mod.}p-1)$ (Artikel 19). Soll nun eine beliebige primitive Wurzel $r$ als Produkt $ABC$ usw. dargestellt werden, so nehme man $A\equiv r^{\mathfrak{a}},B\equiv r^{\mathfrak{b}},C\equiv r^{\mathfrak{c}}$ usw. Dabei gehört $A$ zum Exponenten $a^{\alpha}$, $B$ zum Exponenten $b^{\beta}$ usw. Das Produkt aller $A,B,C$ usw. ist $\equiv r\,(\operatorname{mod.}p)$. Schließlich erkennt man leicht, dass sich $A,B,C$ usw. nicht anders bestimmen lassen.}.

Diese Produkte können also an die Stelle der primitiven Wurzeln selbst treten. Da in ihnen alle Werte von $A$ mit allen Werten von $B$ usw. zu verbinden sind, ist die Summe aller Produkte gleich dem Produkt aus der Summe aller Werte von $A$, der Summe aller Werte von $B$, der Summe aller Werte von $C$ usw., wie aus der Kombinationslehre bekannt ist. Bezeichnen wir alle Werte von $A$; $B$ usw. mit $A,A',A''$ usw.; $B,B',B''$ usw. usw., so ist die Summe aller primitiven Wurzeln
\[\equiv(A+A'+\text{usw.})(B+B'+\text{usw.})\text{usw.}\]
Nun behaupte ich: Ist der Exponent $\alpha$ $=1$, so ist die Summe $A+A'+A''+\text{usw.}$ $\equiv-1\,(\operatorname{mod.}p)$; ist dagegen $\alpha$ $>1$, so ist diese Summe $\equiv0$. Entsprechendes gilt für die übrigen Exponenten $\beta,\gamma$ usw. Sobald dies bewiesen ist, liegt die Wahrheit unseres Satzes
\SourceLeaf{69}auf der Hand. Ist nämlich $p-1$ durch eine Quadratzahl teilbar, so ist einer der Exponenten $\alpha,\beta,\gamma$ usw. größer als eins. Damit ist einer der Faktoren, deren Produkt der Summe aller primitiven Wurzeln kongruent ist, $\equiv0$, also auch das Produkt selbst. Ist dagegen $p-1$ durch keine Quadratzahl teilbar, so sind alle Exponenten $\alpha,\beta,\gamma$ usw. $=1$. Die Summe aller primitiven Wurzeln ist dann einem Produkt aus ebenso vielen Faktoren, jeder $\equiv-1$, kongruent, wie es Zahlen $a,b,c$ usw. gibt. Sie ist also $\equiv\pm1$, je nachdem, ob deren Anzahl gerade oder ungerade ist. Die Behauptungen beweist man wie folgt.

$1^0$. Ist $\alpha=1$ und gehört die Zahl $A$ zum Exponenten $a$, so sind die übrigen zu diesem Exponenten gehörigen Zahlen $A^2,A^3\GaussFourDots A^{a-1}$. Nun ist
\[1+A+A^2+A^3\GaussFourDots+A^{a-1}\]
die Summe einer vollständigen Periode, also $\equiv0$ (Artikel 79). Somit ist
\[A+A^2+A^3\GaussFourDots+A^{a-1}\equiv-1\]

$2^0$. Ist dagegen $\alpha>1$ und gehört die Zahl $A$ zum Exponenten $a^{\alpha}$, so erhält man die übrigen zu diesem Exponenten gehörigen Zahlen, indem man aus $A^2,A^3,A^4\GaussFourDots A^{a^{\alpha}-1}$ die Potenzen $A^a,A^{2a},A^{3a}$ usw. weglässt; siehe Artikel 53. Ihre Summe ist daher
\[\equiv1+A+A^2\GaussFourDots+A^{a^{\alpha}-1}-(1+A^a+A^{2a}\GaussFourDots+A^{a^{\alpha}-a})\]
\emph{also} der Differenz zweier Perioden kongruent und somit $\equiv0$. \mbox{\emph{Q. E. D.}}

\begin{center}\emph{Moduln, die Potenzen von Primzahlen sind.}\end{center}

\Article{82}
Alles bisher Vorgetragene beruht auf der Annahme, dass der Modul eine Primzahl ist. Es bleibt der Fall, dass eine zusammengesetzte Zahl als Modul genommen wird. Hier treten aber weder so schöne Eigenschaften wie im vorigen Fall hervor, noch braucht man besondere Kunstgriffe, um sie zu finden. Vielmehr lässt sich fast alles allein durch Anwenden der vorigen Grundsätze gewinnen. Jede Einzelheit auszubreiten wäre daher unnötig und ermüdend. Wir erläutern also kurz, was dieser Fall mit dem vorigen gemeinsam hat und was ihm eigen ist.

\SourceLeaf{70}
\Article{83}
Die Sätze der Artikel 45–48 sind bereits allgemein bewiesen. Den Satz aus Artikel 49 muss man jedoch wie folgt ändern:

\emph{Bezeichnet $f$ die Anzahl der zu $m$ teilerfremden und unter $m$ liegenden Zahlen, gilt also $f=\phi m$ (Artikel 38), so ist der Exponent $t$ der niedrigsten Potenz einer gegebenen zu $m$ teilerfremden Zahl $a$, die nach dem Modul $m$ der Eins kongruent ist, entweder $=f$ oder ein echter Teiler dieser Zahl.}

Der Beweis aus Artikel 49 gilt auch hier, wenn man überall $p$ durch $m$, $p-1$ durch $f$ und die Zahlen $1,2,3,\GaussFourDots p-1$ durch die zu $m$ teilerfremden, unter $m$ liegenden Zahlen ersetzt. Wir verweisen den Leser daher dorthin. Die übrigen dort erwähnten Beweise (Artikel 50, 51) lassen sich dagegen nicht ohne erhebliche Umwege auf diesen Fall anwenden. --- Bei den folgenden Sätzen, Artikel 52 und weiter, tritt aber ein großer Unterschied auf zwischen Moduln, die Potenzen von Primzahlen sind, und solchen, die sich durch mehrere Primzahlen teilen lassen. Wir betrachten daher zuerst die Moduln der ersten Art für sich.

\Article{84}
Ist der Modul $m=p^n$, wobei $p$ eine Primzahl bezeichnet, so gilt $f=p^{n-1}(p-1)$ (Artikel 38). Wendet man die Untersuchungen aus Artikel 53, 54 mit den im vorigen Artikel angegebenen Änderungen auf diesen Fall an, so zeigt sich, dass alles dort Bewiesene auch hier gilt, wenn zuvor bewiesen wäre, dass eine Kongruenz der Form $x^t-1\equiv0\,(\operatorname{mod.}p^n)$ nicht mehr als $t$ verschiedene Wurzeln haben kann. Für einen Primzahlmodul haben wir dies aus dem allgemeineren Satz von Artikel 43 abgeleitet. Dieser gilt aber in seinem ganzen Umfang nur für Primzahlmoduln und ist daher hier nicht anwendbar. Trotzdem beweisen wir den Satz für diesen besonderen Fall durch ein eigenes Verfahren. Weiter unten, in Abschnitt VIII, zeigen wir, wie man dasselbe leichter findet.

\Needspace{10\baselineskip}
\Article{85}
Wir wollen folgenden Satz beweisen:

\emph{Ist der größte gemeinsame Teiler der Zahlen $t$ und $p^{n-1}(p-1)$ gleich $e$, so hat die Kongruenz $x^t\equiv1\,(\mathit{mod.}p^n)$ genau $e$ verschiedene Wurzeln.}

Sei $e=kp^{\nu}$, sodass $k$ den Faktor $p$ nicht enthält und somit die Zahl $p-1$
\SourceLeaf{71}teilt. Dann hat die Kongruenz $x^t\equiv1$ nach dem Modul $p$ genau $k$ verschiedene Wurzeln. Bezeichnen wir sie mit $A,B,C$ usw., so muss jede Wurzel derselben Kongruenz nach dem Modul $p^n$ nach $p$ einer der Zahlen $A,B,C$ usw. kongruent sein. Nun zeigen wir, dass die Kongruenz $x^t\equiv1.\,(\operatorname{mod.}p^n)$ genau $p^{\nu}$ Wurzeln hat, die $A$, ebenso viele, die $B$ usw. nach dem Modul $p$ kongruent sind. Damit ist die Anzahl aller Wurzeln $kp^{\nu}$ oder $e$, wie behauptet. Den Beweis führen wir so: \emph{Zuerst} zeigen wir, dass mit einer Wurzel $\alpha$, die $A$ nach dem Modul $p$ kongruent ist, auch

\[\alpha+p^{n-\nu},\ \alpha+2p^{n-\nu},\ \alpha+3p^{n-\nu},\mathinner{.\,.\,.\,.}\alpha+(p^{\nu}-1)p^{n-\nu}\]

Wurzeln sind. \emph{Zweitens} können unter den Zahlen, die $A$ nach dem Modul $p$ kongruent sind, nur die von der Form $\alpha+hp^{n-\nu}$ Wurzeln sein, wobei $h$ eine beliebige ganze Zahl bezeichnet. Damit gibt es offenbar $p^{\nu}$ verschiedene Wurzeln und nicht mehr. Dasselbe gilt für die Wurzeln, die den einzelnen Zahlen $B,C$ usw. kongruent sind. \emph{Drittens} zeigen wir, wie sich stets eine Wurzel finden lässt, die $A$ nach $p$ kongruent ist.

\Article{86}
\textsc{Satz.} \emph{Ist wie im vorigen Artikel $t$ durch $p^{\nu}$, aber nicht durch $p^{\nu+1}$ teilbar, so gilt}

\[\begin{aligned}
(\alpha+hp^{\mu})^t-\alpha^t&\equiv0\,(mod.p^{\mu+\nu}),\\
&\textit{aber}\ \equiv\alpha^{t-1}hp^{\mu}t\,(mod.p^{\mu+\nu+1})
\end{aligned}\]

Der zweite Teil des Satzes gilt nicht, wenn zugleich $p=2$ und $\mu=1$ ist.

Der Satz ließe sich durch Entwickeln der Binompotenz beweisen, wenn man zeigte, dass alle Glieder nach dem zweiten durch $p^{\mu+\nu+1}$ teilbar sind. Da die Betrachtung der Nenner der Koeffizienten aber einige Umwege erfordert, ziehen wir das folgende Verfahren vor.

Nehmen wir \emph{zunächst} $\mu>1$ und $\nu=1$ an. Wegen

\[x^t-y^t=(x-y)(x^{t-1}+x^{t-2}y+x^{t-3}y^2+\text{usw.}+y^{t-1})\]

gilt dann

\[\begin{aligned}
(\alpha+hp^{\mu})^t-\alpha^t={}&hp^{\mu}((\alpha+hp^{\mu})^{t-1}\\
&+(\alpha+hp^{\mu})^{t-2}\alpha+\text{usw.}+\alpha^{t-1})
\end{aligned}\]

Nun ist $\alpha+hp^{\mu}\equiv\alpha\,(\operatorname{mod.}p^2)$.
\SourceLeaf{72}Daher ist jedes Glied $(\alpha+hp^{\mu})^{t-1}$, $(\alpha+hp^{\mu})^{t-2}\alpha$ usw. $\equiv\alpha^{t-1}(\operatorname{mod.}p^2)$, also die Summe aller Glieder $\equiv t\alpha^{t-1}(\operatorname{mod.}p^2)$ oder von der Form $t\alpha^{t-1}+Vp^2$, wobei $V$ irgendeine Zahl bezeichnet. Damit hat $(\alpha+hp^{\mu})^t-\alpha^t$ die Form

\[\begin{aligned}
&\alpha^{t-1}hp^{\mu}t+Vhp^{\mu+2},\\
&\textit{das heißt,}\ \equiv\alpha^{t-1}hp^{\mu}t(\operatorname{mod.}p^{\mu+2})\quad\text{und}\ \equiv0\,(\operatorname{mod.}p^{\mu+1}).
\end{aligned}\]

\noindent Für diesen Fall ist der Satz also bewiesen.

Wäre der Satz nun für andere Werte von $\nu$ bei weiterhin geltendem $\mu>1$ falsch, so müsste es eine Grenze geben, bis zu der er stets richtig, darüber hinaus aber falsch wäre. Der kleinste Wert von $\nu$, für den er falsch ist, sei $=\varphi$. Dann sieht man leicht: Ist $t$ durch $p^{\varphi-1}$, aber nicht durch $p^{\varphi}$ teilbar, so gilt der Satz noch; ersetzt man dagegen $t$ durch $tp$, so gilt er nicht. Wir haben also

\[\begin{aligned}
(\alpha+hp^{\mu})^t&\equiv\alpha^t+\alpha^{t-1}hp^{\mu}t(\operatorname{mod.}p^{\mu+\varphi})\\
&\text{oder}\ =\alpha^t+\alpha^{t-1}hp^{\mu}t+up^{\mu+\varphi}
\end{aligned}\]

\noindent wobei $u$ eine ganze Zahl bezeichnet. Da aber der Satz für $\nu=1$ bereits bewiesen ist, gilt

\[\begin{aligned}
(\alpha^t+\alpha^{t-1}hp^{\mu}t+up^{\mu+\varphi})^p\equiv{}&\alpha^{tp}+\alpha^{tp-1}hp^{\mu+1}t\\
&+\alpha^{tp-t}up^{\mu+\varphi+1}\,(\operatorname{mod.}p^{\mu+\varphi+1})
\end{aligned}\]

\noindent und somit auch

\[(\alpha+hp^{\mu})^{tp}\equiv\alpha^{tp}+\alpha^{tp-1}hp^{\mu}tp\,(\operatorname{mod.}p^{\mu+\varphi+1})\]

\noindent \emph{das heißt,} der Satz gilt auch, wenn $t$ durch $tp$ ersetzt wird, \emph{also} auch für $\nu=\varphi$, entgegen der Voraussetzung. Damit gilt der Satz offenbar für sämtliche Werte von $\nu$.

\Article{87}
Es bleibt der Fall $\mu=1$. Mit einem ganz ähnlichen Verfahren wie im vorigen Artikel lässt sich ohne den binomischen Satz zeigen, dass

\[\begin{aligned}
(\alpha+hp)^{t-1}&\equiv\alpha^{t-1}+\alpha^{t-2}(t-1)hp\quad(\operatorname{mod.}p^2)\\
\alpha(\alpha+hp)^{t-2}&\equiv\alpha^{t-1}+\alpha^{t-2}(t-2)hp\\
\alpha\alpha(\alpha+hp)^{t-3}&\equiv\alpha^{t-1}+\alpha^{t-2}(t-3)hp\\
\text{usw.}\quad&
\end{aligned}\]

\noindent also die Summe (da die Anzahl der Glieder $=t$)

\SourceLeaf{73}
\[\equiv t\alpha^{t-1}+\tfrac{(t-1)t}{2}\alpha^{t-2}hp(\operatorname{mod.}p^2)\]

ist. Da $t$ durch $p$ teilbar ist, ist auch $\dfrac{(t-1)t}{2}$ durch $p$ teilbar, außer wenn $p=2$ ist; darauf haben wir schon im vorigen Artikel hingewiesen. In den übrigen Fällen gilt $\tfrac{(t-1)t}{2}\alpha^{t-2}hp\equiv0\,(\operatorname{mod.}p^2)$, und somit ist auch jene Summe $\equiv t\alpha^{t-1}(\operatorname{mod.}p^2)$, wie im vorigen Artikel. Im Übrigen verläuft der Beweis hier genauso wie dort.

Wir schließen also allgemein, mit der einzigen Ausnahme $p=2$, dass

\[(\alpha+hp^{\mu})^t\equiv\alpha^t(\operatorname{mod.}p^{\mu+\nu})\]

\noindent und $(\alpha+hp^{\mu})^t\ \textit{nicht}\equiv\alpha^t$ für jeden Modul, der eine höhere Potenz von $p$ als $p^{\mu+\nu}$ ist, sofern $h$ nicht durch $p$ teilbar ist und $p^{\nu}$ die höchste Potenz von $p$ ist, die $t$ teilt.

\Needspace{8\baselineskip}
Daraus ergeben sich unmittelbar die erste und zweite Behauptung, die wir in Artikel 85 beweisen wollten:

\emph{erstens,} ist $\alpha^t\equiv1$, so auch $(\alpha+hp^{n-\nu})^t\equiv1\,(\operatorname{mod.}p^n)$;

\emph{zweitens,} würde eine Zahl $\alpha'$, die $A$ und damit auch $\alpha$ nach dem Modul $p$ kongruent ist, nicht aber nach $p^{n-\nu}$, die Kongruenz $x^t\equiv1\,(\operatorname{mod.}p^n)$ erfüllen, so setze man $\alpha'$ $=\alpha+lp^{\lambda}$, wobei $l$ nicht durch $p$ teilbar sei. Dann ist $\lambda<n-\nu$. Damit ist $(\alpha+lp^{\lambda})^t$ nach dem Modul $p^{\lambda+\nu}$ der Zahl $\alpha^t$ kongruent, nicht aber nach dem Modul $p^n$, der eine höhere Potenz ist. Also kann $\alpha'$ keine Wurzel der Kongruenz $x^t\equiv1$ sein.

\Article{88}
Als \emph{Drittes} wollten wir eine Wurzel der Kongruenz $x^t\equiv1\,(\operatorname{mod.}p^n)$ finden, die $A$ kongruent ist. Wir zeigen hier nur, wie dies geht, wenn schon eine Wurzel derselben Kongruenz nach dem Modul $p^{n-1}$ bekannt ist. Das genügt offenbar, da wir vom Modul $p$, für den $A$ eine Wurzel ist, zum Modul $p^2$ und so weiter zu allen aufeinanderfolgenden Potenzen übergehen können.

Sei also $\alpha$ eine Wurzel der Kongruenz $x^t\equiv1\,(\operatorname{mod.}p^{n-1})$, und gesucht sei eine Wurzel derselben Kongruenz nach dem Modul $p^n$. Setzen wir diese $=\alpha+hp^{n-\nu-1}$, denn dass sie diese Form haben muss, folgt aus dem vorigen Artikel. (Den Fall $\nu=n-1$
\SourceLeaf{74}betrachten wir später gesondert; größer als $n-1$ kann $\nu$ aber nicht sein.) Es muss daher gelten

\[(\alpha+hp^{n-\nu-1})^t\equiv1(\operatorname{mod.}p^{n-1})\]

\noindent Aber\hfill $(\alpha+hp^{n-\nu-1})^t\equiv\alpha^t+\alpha^{t-1}htp^{n-\nu-1}\,(\operatorname{mod.}p^n)$\hfill\mbox{}

\noindent Bestimmt man also $h$ so, dass $1\equiv\alpha^t+\alpha^{t-1}htp^{n-\nu-1}\,(\operatorname{mod.}p^n)$ gilt, oder so, dass (da nach Voraussetzung $1\equiv\alpha^t(\operatorname{mod.}p^{n-1})$ und $t$ durch $p^{\nu}$ teilbar ist) $\dfrac{\alpha^t-1}{p^{n-1}}+\alpha^{t-1}h\dfrac{t}{p^{\nu}}$ durch $p$ teilbar wird, so ist das Verlangte erreicht. Dass dies immer möglich ist, folgt aus dem vorigen Abschnitt. Denn hier setzen wir voraus, dass $t$ durch keine höhere Potenz von $p$ als $p^{\nu}$ teilbar ist, also $\alpha^{t-1}\dfrac{t}{p^{\nu}}$ zu $p$ teilerfremd ist.

Ist dagegen $\nu=n-1$, \emph{das heißt} $t$ durch $p^{n-1}$ oder sogar durch eine höhere Potenz von $p$ teilbar, so erfüllt jeder Wert $A$, der der Kongruenz $x^t\equiv1$ nach dem Modul $p$ genügt, diese auch nach dem Modul $p^n$. Sei nämlich $t=p^{n-1}\tau$. Dann gilt $t\equiv\tau(\operatorname{mod.}p-1)$. Wegen $A^t\equiv1\,(\operatorname{mod.}p)$ gilt also auch $A^{\tau}\equiv1\,(\operatorname{mod.}p)$. Setzen wir daher $A^{\tau}=1+hp$, so ist $A^t=(1+hp)^{p^{n-1}}\equiv1(\operatorname{mod.}p^n)$ nach Artikel 87.

\Article{89}
Alles, was wir in Artikel 57 und den folgenden mithilfe des Satzes hergeleitet haben, dass die Kongruenz $x^t\equiv1$ nicht mehr als $t$ verschiedene Wurzeln hat, gilt auch für einen Modul, der eine Potenz einer Primzahl ist. Nennen wir die zum Exponenten $p^{n-1}(p-1)$ gehörigen Zahlen \emph{primitive Wurzeln}, also die Zahlen, deren Perioden alle nicht durch $p$ teilbaren Zahlen enthalten, so gibt es auch hier primitive Wurzeln. Alles oben über Indizes und ihre Verwendung sowie über die Lösung der Kongruenz $x^t\equiv1$ Gesagte ist auch hier anwendbar. Da dies keine Schwierigkeiten bietet, wäre es unnötig, alles von Anfang an zu wiederholen. Außerdem haben wir gezeigt, wie sich die Wurzeln der Kongruenz $x^t\equiv1$ nach dem Modul $p^n$ aus den Wurzeln derselben Kongruenz nach $p$ ableiten lassen. Zum Fall, dass eine Potenz von 2 der Modul ist, müssen wir aber noch einiges hinzufügen, da er oben ausgenommen war.

\Needspace{12\baselineskip}
\begin{center}\emph{Moduln, die Potenzen von zwei sind.}\end{center}

\Article{90}
\emph{Wird eine höhere als die zweite Potenz von 2, etwa $2^n$, als Modul genommen, so ist die Potenz mit dem Exponenten $2^{n-2}$ einer beliebigen ungeraden Zahl der Eins kongruent.}
\SourceLeaf{75}\emph{Zum Beispiel} $3^8=6561\equiv1\,(\operatorname{mod.}32)$.

Jede ungerade Zahl hat nämlich entweder die Form $1+4h$ oder $-1+4h$. Daraus folgt der Satz unmittelbar (Satz aus Artikel 86).

Da also der Exponent, zu dem eine beliebige ungerade Zahl nach dem Modul $2^n$ gehört, ein Teiler von $2^{n-2}$ sein muss, gehört sie zu einer der Zahlen $1,2,4,8\mathinner{.\,.\,.\,.}2^{2n-2}$. Zu welcher, lässt sich wie folgt leicht entscheiden. Die gegebene Zahl sei $=4h\pm1$, und der Exponent der höchsten Potenz von 2, die $h$ teilt, sei $=m$ (er kann auch $=0$ sein, nämlich wenn $h$ ungerade ist). Dann ist der Exponent, zu dem die gegebene Zahl gehört, $=2^{n-m-2}$, sofern $n>m+2$ ist. Ist dagegen $n=$ oder $<m+2$, so ist die gegebene Zahl $\equiv\pm1$ und gehört daher entweder zum Exponenten 1 oder zum Exponenten 2. Dass eine Zahl der Form $\pm1+2^{m+2}k$ (gleichbedeutend mit $4h\pm1$), zur Potenz mit dem Exponenten $2^{n-m-2}$ erhoben, der Eins nach dem Modul $2^n$ kongruent wird, bei einem Exponenten, der eine niedrigere Potenz von 2 ist, dagegen nicht, ergibt sich ohne Schwierigkeit aus Artikel 86. Jede Zahl der Form $8k+3$ oder $8k+5$ gehört also zum Exponenten $2^{n-2}$.

\Article{91}
Hier gibt es somit keine \emph{primitiven Wurzeln} in dem oben verwendeten Sinn: keine Zahlen, deren Periode sämtliche unter dem Modul liegenden und zu ihm teilerfremden Zahlen umfasst. Trotzdem sieht man leicht etwas Entsprechendes. Eine Potenz mit ungeradem Exponenten einer Zahl der Form $8k+3$ hat stets die Form $8k+3$, eine mit geradem Exponenten dagegen stets die Form $8k+1$. Keine Potenz kann also die Form $8k+5$ oder $8k+7$ haben. Nun besteht die Periode einer Zahl der Form $8k+3$ aus $2^{n-2}$ verschiedenen Gliedern, jedes von der Form $8k+3$ oder $8k+1$, und es gibt nicht mehr als $2^{n-2}$ solche unter dem Modul liegenden Zahlen. Daher ist jede Zahl der Form $8k+1$ oder $8k+3$ nach dem Modul $2^n$ irgendeiner Potenz einer beliebigen Zahl der Form $8k+3$ kongruent. Ebenso zeigt man, dass die Periode einer Zahl der Form $8k+5$ alle Zahlen der Formen $8k+1$ und $8k+5$ umfasst. Wählt man also eine Zahl der Form $8k+5$ als Basis, so erhalten alle Zahlen der Formen $8k+1$ und $8k+5$, positiv genommen, sowie alle Zahlen der Formen $8k+3$ und $8k+7$, negativ genommen, wirkliche Indizes. Dabei sind nach $2^{n-2}$ kongruente Indizes als gleichwertig anzusehen. So ist unsere Tafel I zu verstehen, in der wir für die Moduln 16, 32 und 64
\SourceLeaf{76}(für den Modul $8$ braucht man keine Tafel) stets die Zahl $5$ als Basis gewählt haben. \emph{Zum Beispiel} entspricht der Zahl $19$, die von der Form $8n+3$ und daher \emph{negativ} zu nehmen ist, beim Modul $64$ der Index $7$. Das bedeutet $5^7\equiv-19\,(\operatorname{mod.}64)$.

Den Zahlen der Formen $8n+1$, $8n+5$, negativ genommen, sowie denen der Formen $8n+3$, $8n+7$, positiv genommen, müsste man dagegen gewissermaßen imaginäre Indizes zuordnen. Führt man sie ein, so lässt sich die Rechnung mit Indizes auf einen sehr einfachen Algorithmus zurückführen. Eine völlig strenge Darstellung würde uns jedoch zu weit führen. Wir behalten uns dies daher für eine andere Gelegenheit vor, wenn wir vielleicht die Theorie der imaginären Größen ausführlicher behandeln werden. Unserer Ansicht nach hat sie bisher niemand auf klare Begriffe zurückgeführt. Kundige werden diesen Algorithmus leicht selbst finden; weniger Geübte können die Tafel trotzdem verwenden, so wie auch diejenigen Logarithmen benutzen, die die Erörterungen neuerer Autoren über imaginäre \emph{Logarithmen} nicht kennen, sofern sie nur die oben begründeten Prinzipien gut verstanden haben.

\begin{center}\small\emph{Aus mehreren Primzahlen zusammengesetzte Moduln.}\end{center}

\Article{92}
Bei einem aus mehreren Primzahlen zusammengesetzten Modul lässt sich fast alles über die Reste von Potenzen aus der allgemeinen Kongruenztheorie ableiten. Da wir aber weiter unten ausführlicher zeigen, wie man beliebige Kongruenzen nach einem solchen Modul auf Kongruenzen mit Primzahlen oder Primzahlpotenzen als Modul zurückführt, brauchen wir uns hier nicht lange damit aufzuhalten. Wir bemerken nur, dass die sehr schöne Eigenschaft der übrigen Moduln, nämlich dass es stets Zahlen gibt, deren Periode sämtliche zum Modul teilerfremden Zahlen umfasst, hier fehlt. Es gibt nur eine Ausnahme: wenn der Modul das Doppelte einer Primzahl oder einer Primzahlpotenz ist.
Bringt man nämlich den Modul $m$ auf die Form $A^aB^bC^c$ usw., wobei $A,B,C$ usw. verschiedene Primzahlen bezeichnen, und bezeichnet man außerdem $A^{a-1}(A-1)$ mit $\alpha$, $B^{b-1}(B-1)$ mit $\beta$ usw., und ist schließlich $z$ eine zu $m$ teilerfremde Zahl, so gilt $z^{\alpha}\equiv1\,(\operatorname{mod.}A^a)$. $z^{\beta}\equiv1\,(\operatorname{mod.}B^b)$ usw. Ist also $\mu$ das kleinste gemeinsame Vielfache der Zahlen $\alpha,\beta,\gamma$ usw., so gilt $z^{\mu}\equiv1$ nach sämtlichen Moduln $A^a,B^b$ usw. und damit auch nach $m$, das deren Produkt ist. Außer wenn $m$ das Doppelte einer Primzahl oder einer Primzahlpotenz ist, ist aber das kleinste gemeinsame Vielfache der Zahlen $\alpha,\beta,\gamma$ usw. kleiner als ihr Produkt. Denn die Zahlen $\alpha,\beta,\gamma$ usw. können nicht paarweise teilerfremd sein, sondern haben jedenfalls den gemeinsamen Teiler $2$. Also kann die Periode keiner Zahl
\SourceLeaf{77}so viele Glieder umfassen, wie es unter dem Modul liegende, zu ihm teilerfremde Zahlen gibt, da deren Anzahl gleich dem Produkt von $\alpha,\beta,\gamma$ usw. ist. So ist \emph{zum Beispiel} für $m=1001$ jede zu $m$ teilerfremde Zahl, zur Potenz mit dem Exponenten $60$ erhoben, der Eins kongruent, da $60$ ein gemeinsames Vielfaches der Zahlen $6,10,12$ ist. --- Der Fall, dass der Modul das Doppelte einer Primzahl oder einer Primzahlpotenz ist, entspricht aber vollständig dem Fall einer Primzahl oder einer Primzahlpotenz als Modul.

\Needspace{17\baselineskip}
\Article{93}
Die Schriften, in denen andere Mathematiker den Gegenstand dieses Abschnitts behandelt haben, wurden bereits an verschiedenen Stellen erwähnt. Wer aber manches ausführlicher erläutert sehen möchte, als es uns die Kürze erlaubte, sei besonders auf die folgenden Abhandlungen Eulers verwiesen. Sie sind wegen der Klarheit besonders zu empfehlen, durch die sich dieser große Mann stets vor allen anderen auszeichnete.

\emph{Theoremata circa residua ex divisione potestatum relicta Comm. nov. Petr. T.} VII, \emph{p.} 49 und folgende.

\emph{Demonstrationes circa residua ex divisione potestatum per numeros primos resultantia.}\par
\hspace*{\parindent}\emph{Ibid. T.} XVIII, \emph{p.} 85 und folgende.

Hinzu kommen können \emph{Opusculorum analyt. T.} I, Abhandlungen 5 und 8.

\begin{center}\rule{.28\linewidth}{.25pt}\end{center}
